\section{Introdução}
A existência de um jogo implementado não demonstra, por si só, que todas as decisões de design foram preservadas ou que os discentes desenvolveram competências de programação; essas dimensões precisam de evidências distintas \citep{prather2024}. Neste exemplo, o problema é identificar a relação entre uma especificação textual e as escolhas necessárias para torná-la executável.

A pergunta demonstrada é: como decisões explícitas do GDD se relacionam com o código e quais lacunas impedem uma comparação direta? O objetivo é organizar essa comparação em registros auditáveis, delimitando a unidade de análise à decisão documentada.

O modelo ilustra uma estrutura de artigo, sem impor capa, folha de aprovação, sumário, índice ou lombada de um trabalho monográfico; a diferenciação entre gêneros é sustentada pelas normas identificadas nas fontes de normalização e pelo exemplo canônico de artigo do abnTeX2 \citep{ufes2026,abntex2018}.

\section{Fundamentação}
Prather et al. analisam benefícios e danos de IA generativa para iniciantes, oferecendo fundamento para distinguir desempenho assistido e competência individual; o presente exemplo utiliza essa cautela para delimitar suas conclusões \citep{prather2024}. O GDD e o código permitem examinar decisões técnicas, mas não constituem uma medida de aprendizagem por definição \citep{prather2024}.

A base de apresentação utiliza requisitos efetivamente lidos nos PDFs do projeto, com atualização das edições identificada em fonte universitária; versões antigas de templates não são tomadas como prova automática de conformidade atual \citep{ufes2026}.
